\documentclass[11pt,a4paper]{article}
\usepackage[margin=1in]{geometry}
\usepackage{amsmath,amssymb}
\usepackage{parskip}
\usepackage{hyperref}

\title{\textbf{Pre-Experiment Overview:\\ Photoelectric Effect}}
\author{IBRAHIM H.I. ABUSHAWISH \\ Istanbul University, Department of Physics}
\date{}

\begin{document}
\maketitle

\section*{Purpose of the Experiment}
According to the lab manual, the purpose of this experiment is to observe the photoelectric phenomenon experimentally, demonstrate the failure of classical wave theory, and determine Planck's constant, $h$, along with the photocathode's work function, $\Phi$, by measuring the stopping voltage as a function of the frequency of incident light.

\section*{Importance and Significance}
The photoelectric effect is historically one of the most decisive milestones in the birth of quantum mechanics. Classical electromagnetic wave theory predicted that electron ejection energy should depend on light intensity and that faint light would require a measurable time delay to liberate electrons. Instead, experiments by Lenard and Millikan proved that electron emission occurs instantaneously, and the maximum kinetic energy depends solely on light frequency, regardless of intensity. Albert Einstein explained this in 1905 by proposing that light behaves as localized packets of energy---photons---each carrying energy $E = h \nu$.

In the experimental setup, light from a high-pressure mercury lamp is passed through optical color filters to select discrete spectral wavelengths (e.g., yellow, green, blue, and violet). The monochromatic light strikes the photocathode in a phototube, liberating photoelectrons. By applying an opposing variable retarding potential (stopping voltage, $V_0$) across the collector plate until the photocurrent drops to zero, we measure the maximum kinetic energy of the ejected electrons: $e V_0 = E_{k,\text{max}} = h \nu - \Phi$. Plotting stopping voltage $V_0$ versus light frequency $\nu$ yields a linear graph whose slope is $h / e$ and whose vertical intercept gives the work function $\Phi$. Calculating Planck's constant from this slope provides direct confirmation of light quantization.

\end{document}
